% Original 15671 architecture: independent CM branch, version 1.
% This lemma is conditional on the separately audited field/ramification lemma.
\paragraph{The CM height estimate over the same field}
Let $A/L$ be an abelian surface of CM type, let $B$ be a quaternion
division algebra over $\mathbb Q$ embedded in
$\operatorname{End}^0_{\overline L}(A)$, and suppose that all geometric
endomorphisms of $A$ are defined over $L$. Assume
\[
 d=[L:\mathbb Q]\le576,\qquad
 \operatorname{rad}(|D_L|)\le R,\qquad R\ge1.
\]
All heights below are stable Faltings heights in Faltings's original
normalization. We prove
\begin{equation}\label{eq:original-cm-bound}
 h_{\rm F}(A)<854R<6^{100}R^{24}.
\end{equation}
No assertion about the existence or degree of $L$ is included in this lemma.

First, $A_{\overline L}$ is isogenous to $E^2$ for a CM elliptic curve
$E$. Here is a dimension argument which makes the alternative cases
explicit. If $A$ were geometrically simple, its division endomorphism
algebra would act faithfully on the four-dimensional rational homology.
Its dimension therefore divides four. As it contains $B$, it would equal
$B$; but $B$ has no commutative semisimple subalgebra of dimension four,
contrary to CM type. If the two elliptic factors were not isogenous,
their endomorphism algebra would be a product of commutative fields and
could not contain $B$. Finally, if $E$ had no CM, the endomorphism algebra
would be $M_2(\mathbb Q)$, and the embedding of the four-dimensional
division algebra $B$ would be an impossible isomorphism.

The elliptic factors and the isogeny may be taken over $L$. Indeed,
the rational idempotents in $\operatorname{End}^0(A)$ are defined over
$L$, and the image of an integral multiple of such an idempotent is an
elliptic subvariety over $L$. Endomorphisms and homomorphisms of the
resulting factors extend to rational endomorphisms of $A$ by the
inclusions and rational projections; these are fixed by
$\operatorname{Gal}(\overline L/L)$. Thus the factor isogenies, and all
endomorphisms of $E$, are defined over $L$.

Put $F=\operatorname{End}^0(E)$, an imaginary quadratic field, and
$\mathcal O=\operatorname{End}(E)$. Differentiation gives an injective
unital map $F\longrightarrow\operatorname{End}_L(\operatorname{Lie}E)=L$:
it is nonzero, and every nonzero elliptic endomorphism in characteristic
zero is separable. Consequently $F\subset L$.

We next pass to the maximal order without extending $L$. Choose
$f\ge1$ with $f\mathcal O_F\subset\mathcal O$, choose an integral
basis $\alpha_1,\alpha_2$ of $\mathcal O_F$, and form the finite
subgroup scheme
\[
 K=\sum_{i=1}^2(f\alpha_i)(E[f])\subset E.
\]
Each $f\alpha_i$ is defined over $L$, so $K$ and the quotient
$E'=E/K$ are defined over $L$. Under a complex uniformization
$E(\mathbb C)=\mathbb C/\Lambda$, this subgroup is
$\mathcal O_F\Lambda/\Lambda$; thus $E'(\mathbb C)$ has lattice
$\mathcal O_F\Lambda$ and endomorphism ring $\mathcal O_F$.
The ring cannot be larger, since the rational endomorphism algebra
remains $F$. This is also a direct special case of the same-field
maximal-order theorem of R\'emond, but the construction just given is
the proof used here. It follows that $A'=(E')^2$ is $L$-isogenous to $A$.

Write $\Delta_F$ for the fundamental discriminant of $F$. The
discriminant tower formula implies that every rational prime dividing
$\Delta_F$ divides $D_L$. A quadratic fundamental discriminant has
odd prime exponents one and exponent at two at most three. Therefore
\[
 |\Delta_F|\le4\operatorname{rad}(|D_L|)\le4R.
\]
The multiplier four includes the possible factor $2^3/2$.

For completeness we give an elementary analytic bound on $h_{\rm F}(E')$.
CM elliptic curves have potentially good reduction. After a finite
extension with everywhere good reduction, the elliptic discriminant
formula of Silverman, reproduced as equation (6.2) in
\cite{cmheight}, gives
\[
 h_{\rm F}(E')=\frac1{[L':\mathbb Q]}\sum_{\sigma:L'\hookrightarrow\mathbb C}
 \left(-\log(2\pi)-\log|\eta(\tau_\sigma)|^2
                     -\tfrac12\log\operatorname{Im}\tau_\sigma\right).
\]
Here $L'$ is used only to evaluate the stable height; it is not the
field used in the isogeny estimate below. The discriminant term is zero
because reduction is good. We use the individual-embedding formula
(6.2), not the class-average assertion of Proposition 6.1, so no
hypothesis on Weierstrass coefficients lying in $\mathbb Q(j)$ is needed.

For each embedding choose a reduced quadratic parameter of discriminant
$\Delta_F$. If $y=\operatorname{Im}\tau$, then
\[
 \frac{\sqrt3}{2}\le y=\frac{\sqrt{|\Delta_F|}}{2a}
                   \le\frac{\sqrt{|\Delta_F|}}2,
 \qquad a\in\mathbb Z_{\ge1}.
\]
Put $r=|e^{2\pi i\tau}|<1/16$. The product defining $\eta$ gives
\[
 -\log|\eta(\tau)|^2
 \le\frac{\pi y}{6}+2\sum_{n\ge1}\frac{r^n}{1-r^n}
 \le\frac{\pi y}{6}+\frac{2r}{(1-r)^2}
 <\frac{\pi\sqrt{|\Delta_F|}}{12}+\frac{32}{225}.
\]
Also $-\tfrac12\log y\le\tfrac14\log(4/3)<1/12$.
Discarding the negative term $-\log(2\pi)$ and using $\pi<4$ shows
\[
 h_{\rm F}(E')<\sqrt{|\Delta_F|}+1\le2\sqrt R+1\le3R.
\]
The height is additive on products, hence $h_{\rm F}(A')<6R$.

Apply \cite[Theorem 1.9(1)]{GaudronRemond2023} over $L$, with source
$A'$ and target $A$ and with no extension of the base field.
For a surface, its quantity $\Xi(A')$ is
\[
 \Xi(A')=\left(14^{32}d\max(1,\log d,h_{\rm F}(A'))\right)^8.
\]
The general assertion supplies an isogeny $A'\to A$ of degree at most
$\Xi(A')^2$; no identification of CM type with that paper's
separate ``MM'' condition is made or needed. The isogeny-height
inequality then gives
\[
 h_{\rm F}(A)\le6R+
 8\log\left(14^{32}\,576\max(1,\log576,6R)\right).
\]
One can see the last inequality directly from invariant differentials:
after semistable base change, pullback on determinant differentials has
integral cokernel, while its archimedean squared norm is multiplied by
the isogeny degree. Thus the height difference is half the logarithm
of that degree minus a nonnegative normalized cokernel length.

Finally, $\log14<3$, $\log576<7$, and
$\max(1,\log576,6R)\le7R$. Since $\log7<2$ and $\log R\le R$,
\[
 h_{\rm F}(A)<6R+8(96+7+2+\log R)
             \le854R.
\]
As $6^4=1296>854$ and $R\ge1$, this proves
\eqref{eq:original-cm-bound}. In particular the original
$6^{100}R^{24}$ allowance absorbs this CM estimate with the field
degree still exactly bounded by 576.
